\documentclass[12pt]{article}
\usepackage[T1]{fontenc}
\usepackage[utf8]{inputenc}
\usepackage{lmodern}
\usepackage{textcomp}
\usepackage{enumitem}
\usepackage{geometry}
\geometry{margin=1in}
\begin{document}
\begin{center}
Answer Key - Cybersecurity - Undergraduate Level Assessment 19 \
\small (Answers are ordered exactly as in the question paper.)
\end{center}

\section*{Multiple Choice}
\begin{enumerate}[label=\arabic*.]
\item \textbf{Answer:} A
\\ \textit{Options:} A) Restrict what operations/data the user can access; B) Determine if the user is an attacker; C) Flag the user if he/she misbehaves; D) Determine who the user is
\\ \textit{Note:} Correct option: A - Restrict what operations/data the user can access
\item \textbf{Answer:} C
\\ \textit{Options:} A) boundary hacks; B) memory checks; C) boundary checks; D) buffer checks
\\ \textit{Note:} Correct option: C - boundary checks
\item \textbf{Answer:} B
\\ \textit{Options:} A) Memory leakage; B) Buffer-overrun; C) Less processing power; D) Inefficient programming
\\ \textit{Note:} Correct option: B - Buffer-overrun
\item \textbf{Answer:} A
\\ \textit{Options:} A) It generates each different input by modifying a prior input; B) It works by making small mutations to the target program to induce faults; C) Each input is mutation that follows a given grammar; D) It only makes sense for file-based fuzzing, not network-based fuzzing
\\ \textit{Note:} Correct option: A - It generates each different input by modifying a prior input
\item \textbf{Answer:} B
\\ \textit{Options:} A) Confidentiality; B) Availability; C) Correctness; D) Integrity
\\ \textit{Note:} Correct option: B - Availability
\end{enumerate}

\section*{True / False}
\begin{enumerate}[label=\arabic*.]
\setcounter{enumi}{5}
\item \textbf{Answer:} True
\\ \textit{Options:} A) True; B) False
\\ \textit{Note:} LLM-generated true statement based on MCQ.
\item \textbf{Answer:} False
\\ \textit{Options:} A) True; B) False
\\ \textit{Note:} LLM-generated false statement. Correct answer: 'buffer'.
\item \textbf{Answer:} False
\\ \textit{Options:} A) True; B) False
\\ \textit{Note:} LLM-generated false statement. Correct answer: 'Whitebox'.
\item \textbf{Answer:} True
\\ \textit{Options:} A) True; B) False
\\ \textit{Note:} LLM-generated true statement based on MCQ.
\item \textbf{Answer:} False
\\ \textit{Options:} A) True; B) False
\\ \textit{Note:} LLM-generated false statement. Correct answer: 'IM - Trojans'.
\end{enumerate}

\section*{Long Form Response}
\begin{enumerate}[label=\arabic*.]
\setcounter{enumi}{10}
\item \textbf{Answer:} def is\_majority(arr, n, x): | return False | return True | return False | def binary\_search(arr, low, high, x): | return mid | return binary\_search(arr, (mid + 1), high, x) | return binary\_search(arr, low, (mid -1), x) | return -1
\\ \textit{Note:} Students should provide a detailed explanation referencing algorithm steps.
\item \textbf{Answer:} def are\_Equal(arr 1,arr 2,n,m): | return False | return False | return True
\\ \textit{Note:} Students should provide a detailed explanation referencing algorithm steps.
\end{enumerate}

\vfill
\noindent\textit{End of Answer Key}
\end{document}